\section{Introduction}
\label{sec:introduction}

The three-dimensional pose of a small molecule determines which chemical
groups can contact a protein, which interactions can be formed, and which
modifications are sterically accessible. Reliable complex structures therefore
support the interpretation of structure--activity relationships, selectivity,
and resistance mutations, and provide starting hypotheses for medicinal
chemistry and more expensive physics-based calculations. Experimental complex
structures remain costly and incomplete, however, particularly for new targets
and ligand series. Protein--ligand docking addresses this gap by predicting a
ligand pose conditioned on a receptor structure and, in the setting considered
here, a specified binding site. This is a geometric prediction problem:
recovering a plausible pose is important for structure-based design, but a
docking score or pose-confidence score is not by itself a binding-affinity
measurement.

Docking is difficult because ligand placement and internal conformation are
strongly coupled. A candidate must resolve global translation and orientation,
multiple rotatable bonds, local protein--ligand complementarity, and
intramolecular geometry on a rugged, multimodal landscape. Protein motion,
uncertain protonation and tautomeric states, cofactors, metals, and solvent can
further change the relevant interaction pattern. Classical docking programs
combine a conformational search with an approximate scoring function and have
become indispensable, practical tools \citep{trott2010vina}; nevertheless,
tractable search often requires a rigid or only locally flexible receptor.
Ligand-induced conformational changes can consequently create a substantial
gap between docking into an available apo structure and redocking into the
experimental holo structure \citep{lu2024dynamicbind}.

Deep generative models have recently changed the formulation of the problem.
DiffDock casts ligand docking as diffusion over translational, rotational, and
torsional degrees of freedom and samples multiple candidate poses
\citep{corso2023diffdock}. DynamicBind extends geometric generation toward
ligand-specific receptor conformations \citep{lu2024dynamicbind}, while
AlphaFold~3 demonstrates that raw all-atom coordinates of proteins, ligands,
nucleic acids, ions, and modified residues can be generated in a unified
cofolding framework \citep{abramson2024alphafold3}. In parallel, FlowDock
applies conditional geometric flow matching to flexible protein--ligand
complex generation \citep{morehead2025flowdock}. Collectively, these results
show that learned generative vector fields can explore complex pose
distributions without relying exclusively on hand-designed search heuristics.
They also make the distinction between tasks more important: blind cofolding
from sequence, docking into a predicted apo receptor, known-pocket docking,
and reference-pocket holo redocking expose models to different information and
should not be compared as if they were the same prediction problem.

Despite this progress, three limitations remain central. First, sampling and
selection are different problems. A generator may place a near-native pose
among many samples while its confidence model fails to rank that pose first.
Second, coordinate accuracy and molecular validity are different axes:
low heavy-atom RMSD does not ensure valid stereochemistry, bond geometry, or
absence of severe clashes \citep{buttenschoen2024posebusters}. Third, random
row-level splits can overstate generalization when related ligands, homologous
pockets, or alternate structures cross the train--test boundary. Recent
systematic studies have found that performance changes materially for predicted
apo receptors, unknown pockets, multiligand systems, and previously unseen
protein--ligand poses \citep{morehead2026posebench}. A post-training-cutoff
study of 2,600 high-resolution systems further reported strong dependence of
cofolding accuracy on similarity to training examples
\citep{skrinjar2026generalization}. These observations motivate
similarity-aware resources such as PLINDER \citep{durairaj2024plinder} and
require every docking result to declare the receptor state, pocket
information, sampling budget, ranking rule, and leakage boundary.

The representation used for ligand flexibility is another consequential
choice. Treating the entire ligand as one rigid body preserves chemistry but
cannot express torsional rearrangement. Directly generating every atom is
maximally flexible, but requires the model to maintain covalent geometry while
solving the global docking problem. Torsion-based models preserve bond lengths
and angles, yet a rotation about one bond moves an entire downstream subgraph
and ties the parameterization to a particular molecular-tree convention. We
instead use rigid fragments as an intermediate geometric scale. Rotatable bonds
partition the ligand into locally rigid components, and the pose is represented
by one translation in \(\R^3\) and one proper rotation in \(\SO(3)\) per
fragment. This factorization preserves internal fragment geometry while
allowing cooperative relative motion across flexible bonds. Because independent
fragment poses can strain the cut-bond interfaces, the representation must be
paired with explicit fragment adjacency, cross-fragment geometry, and
atom-level contact reasoning.

We introduce \method for pocket-conditioned protein--ligand docking. It uses
fragment-based \(\SE(3)\)-equivariant flow matching. Given a receptor
structure, ligand graph, and externally supplied pocket center,
\method crops a \(10\,\angstrom\) receptor neighborhood and constructs a
heterogeneous graph of ligand atoms, ligand fragments, protein atoms, and
residue nodes. A shared equivariant trunk communicates through covalent,
fragment, residue, and dynamically rebuilt protein--ligand contact edges. The
network predicts equivariant atom-level force-like vectors and aggregates them
into fragment translations and angular velocities through a Newton--Euler
construction. Flow matching learns the time-dependent vector field that
transports independently sampled fragment poses toward the bound pose, while a
separately trained confidence model ranks the resulting candidate set. The
receptor is fixed, and the model neither discovers an unknown pocket nor
predicts binding affinity; this deliberately narrow contract permits a precise
assessment of pose generation and selection.

Our main contributions are:
\begin{itemize}
    \item We formulate flexible ligand docking as flow matching on a product of
    fragment rigid-body pose spaces and introduce an atom-to-fragment
    Newton--Euler head that preserves \(\SE(3)\) equivariance while coupling
    local contact evidence to global fragment motion.

    \item We develop a unified heterogeneous interaction network with
    higher-order irreducible representations, time conditioning, explicit
    cross-fragment geometry, and dynamic protein--ligand contacts, together
    with a matched confidence model for top-1 pose selection.

    \item We separate sampling coverage, deployable selection, and physical
    validity in evaluation. On frozen \(10\,\angstrom\) reference-pocket
    redocking manifests for Astex Diverse and PoseBusters v2, we
    compare the first sample, score-based and learned selectors, and an
    explicitly non-deployable best-of-80 oracle under the same candidate sets.

    \item We state the scientific boundary of the current evidence: the
    reported benchmark is a compatibility diagnostic with target-derived
    pocket centers, not a claim of blind docking, affinity prediction, or
    prospective chemical and pocket generalization. This makes the remaining
    strict-split retraining and apo-receptor tests explicit rather than
    conflating them with holo redocking performance.
\end{itemize}
